% ============ Solutions and hints for exercises (Chapter 7: §19-§21) ============
% 原书 pp.294-295; 原书对 §19 的 Ex.19.1,19.2、§20 的 Ex.20.1,20.2,20.3,20.6 未给解答, 编号按原书跳号.

\section*{\S\ref{sec:19}}

\noindent\textbf{19.3.}\quad If $0 \to P_r \to \cdots \to P_0 \to M \to 0$ is an exact sequence and each $P_i$ is finite
and projective, and if $P_r \oplus A^m = A^n$, then $0 \to P_r \oplus A^m \to P_{r-1} \oplus A^m \to \cdots \to P_0 \oplus A^m \to M \to 0$
is again exact, with $P_r \oplus A^m$ free. Proceeding in the same
way, adding a free module to $P_i$ at each stage, we get an FFR $0 \to L_{r+1}
\to L_r \to \cdots \to L_0 \to M \to 0$.

\noindent\textbf{19.4.}\quad For every maximal ideal $m$ of $A$, since the $A$-module $A/m$ has an FFR, also
the $A_m$-module $A_m/mA_m$ has an FFR, so that the projective dimension over
$A_m$ is finite. Thus $A_m$ is a regular local ring.

\section*{\S\ref{sec:20}}

\noindent\textbf{20.4.}\quad Use Theorem~\hyperref[thm:20.5]{5} and Ex.~\ref{ex:8.3-ex}.

\noindent\textbf{20.5.}\quad This follows from Theorem~\hyperref[thm:20.1]{1}. If an ideal $I$ is locally principal it is finitely
generated, and principal by Theorem~\ref{thm:5.8}.

\section*{\S\ref{sec:21}}

\noindent\textbf{21.1.}\quad For $P \in \Spec R$ with $P \supset I$ set $P/I = p$; then $A_p$ is c.i. $\iff I_p$ is generated by an
$R_p$-sequence $\iff (I/I^2)_p$ is free over $A_p$. Now $I/I^2$ is a finite $A$-module, so that
by Theorem~\ref{thm:4.10}, $\{p \in \Spec A \mid (I/I^2)_p \text{ is free}\}$ is an open subset of $\Spec A$.

\noindent\textbf{21.2.}\quad We can assume that $A$ is complete. Then $A = R/I$ with $R$ regular and
$\dim R = \dim A + 1$. Now $\htop I = 1$ and $A$ is a CM ring, so that all the prime
divisors of $I$ have height 1. Since $R$ is a UFD, $I$ is principal.

\noindent\textbf{21.3.}\quad $A = k[x, y, z] = k + kx + ky + kz + kx^2$, with $x^2 = y^2 = z^2$ and $xy = yz =
zx = 0$. Therefore $0:m = kx^2$, and $A$ is a zero-dimensional Gorenstein
ring. Set $I = (X^2 - Y^2, Y^2 - Z^2, XY, YZ, ZX)$ and $M = (X, Y, Z)$; then
$I \to M^2/M^3$ has five-dimensional image, so that at least five elements
are needed to generate $I$.
